\documentclass[aspectratio=169]{beamer}
\usepackage[T5]{fontenc}
\usepackage[utf8]{inputenc}
\usepackage[vietnamese]{babel}
\usepackage{lmodern}
\usepackage{hyperref}
\usetheme{Madrid}
\title{ConferenceSpace: Tổng quan và bối cảnh đề tài}
\subtitle{Chương 1 -- Mở đầu}
\author{Nhóm thực hiện}
\date{\today}
\begin{document}
\frame{\titlepage}
\begin{frame}{Nội dung}
\tableofcontents
\end{frame}
\section{Đặt vấn đề}
\begin{frame}{Bối cảnh quá tải phản biện}
\begin{itemize}
\item Số lượng bài báo tại các hội nghị AI lớn tăng nhanh trong thập kỷ qua.
\item NeurIPS tăng từ khoảng 9.467 bản thảo năm 2020 lên 21.575 bản thảo năm 2025.
\item Khối lượng cho reviewer, area chair và program chair tăng tương ứng.
\end{itemize}
\end{frame}
\begin{frame}{AI/LLM: cơ hội và rủi ro}
\begin{itemize}
\item LLM có tiềm năng hỗ trợ tóm lược, trích xuất metadata và cải thiện chất lượng phản biện.
\item Rủi ro chính nằm ở bảo mật bản thảo, trách nhiệm học thuật và liêm chính nghiên cứu.
\item AI chỉ nên là công cụ hỗ trợ; quyết định cuối cùng vẫn thuộc về con người.
\end{itemize}
\end{frame}
\begin{frame}{Khoảng trống nền tảng hiện tại}
\begin{itemize}
\item EasyChair, HotCRP, Microsoft CMT và OpenReview hỗ trợ tốt nghiệp vụ quản lý hội nghị truyền thống.
\item Các nền tảng này chưa tập trung vào AI hỗ trợ có kiểm soát cho peer review.
\item ConferenceSpace hướng đến lấp khoảng trống giữa nghiệp vụ, thuật toán xác định và AI hỗ trợ.
\end{itemize}
\end{frame}
\section{Mục tiêu đề tài}
\begin{frame}{Mục tiêu tổng quát}
\begin{itemize}
\item Xây dựng nền tảng web ConferenceSpace phục vụ vòng đời xét duyệt bài báo.
\item Chứng minh mô hình tích hợp AI/LLM minh bạch, có kiểm soát và giải thích được.
\item Giữ nguyên tắc human-in-the-loop trong mọi quyết định học thuật quan trọng.
\end{itemize}
\end{frame}
\begin{frame}{Mục tiêu cụ thể}
\begin{itemize}
\item Quản lý hội nghị, nộp bài, phân công phản biện, đánh giá, rebuttal và ra quyết định.
\item Tách ba lớp xử lý: nghiệp vụ, thuật toán xác định và AI hỗ trợ.
\item Thiết kế AI theo hướng gợi ý, tóm lược, kiểm tra sơ bộ và hỗ trợ tổng hợp thông tin.
\item Đánh giá tính hữu ích, độ tin cậy và mức chấp nhận của người dùng.
\end{itemize}
\end{frame}
\section{Phạm vi}
\begin{frame}{Phạm vi bao gồm}
\begin{itemize}
\item Nghiệp vụ cốt lõi tương đương các hệ thống quản lý hội nghị phổ biến.
\item Reviewer matching dựa trên Domain Jaccard Similarity, độc lập với LLM.
\item Phát hiện COI đa tầng: self-author, khai báo thủ công và đồ thị đồng tác giả.
\item Sáu mô-đun AI hỗ trợ theo vai trò trong quy trình xét duyệt.
\end{itemize}
\end{frame}
\begin{frame}{Phạm vi không bao gồm}
\begin{itemize}
\item Quản lý sự kiện hội nghị như bán vé, xếp lịch phòng hoặc in kỷ yếu.
\item Tự động hóa quyết định chấp nhận hoặc từ chối bài báo.
\item Đánh giá tác động dài hạn của AI lên văn hóa phản biện học thuật.
\item Kiểm soát hoàn toàn chất lượng mô hình LLM bên ngoài.
\end{itemize}
\end{frame}
\section{Cấu trúc luận văn}
\begin{frame}{Cấu trúc sáu chương}
\begin{enumerate}
\item Mở đầu.
\item Khảo sát nhu cầu.
\item Xây dựng hệ thống.
\item Công nghệ sử dụng.
\item Thiết lập thực nghiệm và đánh giá.
\item Kết luận.
\end{enumerate}
\end{frame}
\begin{frame}{Mạch triển khai}
\begin{itemize}
\item Chương 1--2 xác định nhu cầu và khoảng trống.
\item Chương 3--4 cụ thể hóa thiết kế hệ thống và công nghệ.
\item Chương 5 kiểm chứng bằng thực nghiệm và phản hồi người dùng.
\item Chương 6 tổng kết kết quả, hạn chế và hướng phát triển.
\end{itemize}
\end{frame}
\end{document}